\section{Preliminaries}\label{sec:prelim}

This section fixes notation, recalls the prerequisite results we cite (stated, not reproved),
and defines the model. Throughout, $m\in\R$ is a scalar order parameter and $\lambda\ge0$ a
scalar control parameter.

\subsection{Order and control parameters}

\begin{definition}[Order parameter]\label{def:order}
The \emph{misalignment order parameter} $m\in\R$ is the one-dimensional internal-state
coordinate of the system: $m=0$ denotes the aligned state and large $|m|$ denotes a misaligned
state. The physical observable is $|m|$; allowing $m\in\R$ exposes the reflection symmetry that
governs the pitchfork. Empirically, $m$ is identified with a toxic-persona latent activation,
the dominant principal component $|\mathrm{PC1}|$ of trait drift, or a moral-susceptibility
score $S$ \cite{wang2025persona,nghiem2026trait,costa2026collapse}.
\end{definition}

\begin{definition}[Control and bias]\label{def:control}
The \emph{intervention strength} $\lambda\ge0$ is the slow control variable (malicious-data
fraction, training step, or LoRA rank). The \emph{bias field} $h\in\R$ encodes the directional
asymmetry of the intervention. We take the affine calibrations
\begin{equation}\label{eq:calib}
   a(\lambda)=a_0-c\lambda,\qquad h(\lambda)=h_0+e\lambda,
   \qquad a_0,c>0,\ \ e\ge0,
\end{equation}
so that the alignment stiffness $a$ decreases with the intervention.
\end{definition}

\subsection{The free energy and its flow}

\begin{definition}[Effective free energy]\label{def:F}
For a fixed quartic coefficient $b>0$, the \emph{effective free energy} is the cusp normal form
\begin{equation}\label{eq:cusp}
   F(m;a,h)\;=\;\tfrac14\,b\,m^4+\tfrac12\,a\,m^2-h\,m .
\end{equation}
The induced \emph{gradient flow} and its stochastic perturbation are
\begin{equation}\label{eq:flow}
   \dot m=-F'(m)=-(b\,m^3+a\,m-h),\qquad
   dm=-F'(m)\,dt+\sigma\,dW_t,
\end{equation}
with $\sigma>0$ the noise amplitude and $W_t$ a standard Wiener process.
\end{definition}

\begin{definition}[Curvature and stability]\label{def:curv}
The \emph{Hessian} (curvature) of $F$ is $H(m)=F''(m)=3bm^2+a$. An equilibrium $m_e$ of
\eqref{eq:flow} (a root of $F'$) is asymptotically stable if and only if $H(m_e)>0$. We write
$\alpha(\lambda)=H(m_a)$ for the \emph{recovery rate} at the aligned equilibrium $m_a$ and
$\Lambda=-H(m_e)$ for the \emph{Lyapunov exponent} at $m_e$.
\end{definition}

\subsection{Prerequisite results}

We state the results we use; full statements and proofs are in the cited sources. The first is
the variational-flow lemma of the Free Energy Principle.

\begin{theorem}[Free-energy lemma; Friston, Da Costa, Parr~\cite{friston2021}]\label{thm:fep}
Under a Markov-blanket partition at nonequilibrium steady state, the expected flow of the
internal states $\mu$ is $\dot\mu=(Q-\Gamma)\nabla_\mu F$, where $\Gamma$ is a positive-definite
diffusion (mobility) tensor and $Q=-Q^\top$ is solenoidal. The free energy $F$ is a Lyapunov
function for the expected flow and is an upper bound on surprisal $\Im(x)=-\ln p(x)$.
\end{theorem}

\begin{theorem}[Catastrophic-bifurcation classification; Kuehn~\cite{kuehn2011}]\label{thm:kuehn}
A codimension-one equilibrium bifurcation is a critical (catastrophic, jump) transition if and
only if it is a fold/saddle-node, a subcritical Hopf, a subcritical pitchfork, or
transcritical. Supercritical pitchfork and Hopf bifurcations are continuous (non-catastrophic).
\end{theorem}

\begin{theorem}[Gradient dynamic-transition criterion; Ma--Wang~\cite{mawang2008}]\label{thm:mawang}
For a gradient (variational) flow, the dynamic transition at $\lambda_0$ is continuous
(Type-I) if the bifurcating state is locally asymptotically stable at $\lambda_0$, and a jump
(Type-II) otherwise. Equivalently, reducing the flow on the center manifold to
$\dot x=\alpha x^k+o(|x|^k)$ with $k$ odd, the transition is continuous if $\alpha<0$ and a jump
if $\alpha>0$.
\end{theorem}

\begin{theorem}[Early-warning laws; Kuehn~\cite{kuehn2011}]\label{thm:warning}
For fluctuations linearized about an attracting equilibrium with recovery rate $\alpha>0$, the
Ornstein--Uhlenbeck approximation gives stationary variance $\Var(m)\to\sigma^2/(2\alpha)$ and
lag-$\tau$ autocorrelation $\rho(\tau)=e^{-\alpha\tau}$. As $\alpha\to0^+$ on approach to a
transition, both the variance and the autocorrelation time $1/\alpha$ diverge.
\end{theorem}

\begin{theorem}[Noise--timescale scaling; Kuehn~\cite{kuehn2011}]\label{thm:noise}
For a fold with slow-variable speed $\varepsilon$, noise-induced escape before the
deterministic threshold is unlikely if $\sigma\ll\sqrt{\varepsilon}$ and almost sure if
$\sigma\gg\sqrt\varepsilon$, with an intermediate regime at $\sigma\approx\sqrt\varepsilon$
(and $\sigma\approx\varepsilon^{3/4}$ for the pitchfork).
\end{theorem}

We also record, for the limitations discussion, the blanket-degradation result of Aguilera et
al.~\cite{aguilera2022}: in canonical nonequilibrium models the conditional mutual information
$I(x;y\mid b)$ across a Markov blanket grows with entropy production and peaks near the
order--disorder critical point. This bounds the regime in which Theorem~\ref{thm:fep}'s gradient
form is a good approximation.

\subsection{Standing assumptions}

We work under the following assumptions, justified in Section~\ref{sec:discussion}.
\begin{enumerate}[label=(A\arabic*),leftmargin=*,itemsep=2pt,topsep=2pt]
  \item \emph{Quasi-gradient regime.} The solenoidal part $Q$ of Theorem~\ref{thm:fep} is
    neglected; we analyze the symmetric (gradient) part of the flow, which governs the
    stationary density and the transition. The blanket does not dissolve
    (Theorem~\ref{thm:fep} applies; cf.~\cite{aguilera2022}).
  \item \emph{One-dimensional reduction.} The high-dimensional internal state is reduced to the
    scalar $m$, justified by the empirically observed rank-one dominance of the misalignment
    direction \cite{nghiem2026trait,minder2026traces}; the reduction is exact locally on the
    center manifold.
  \item \emph{Positive quartic near threshold.} $b>0$ in a neighborhood of $\lambda^\ast$, so
    $F$ is bounded below and the local expansion \eqref{eq:cusp} is the correct normal form.
\end{enumerate}
